% !PRINTED 560
\weberpage{560}
%<<CONT
then the determinant $D$ of the quadratic form (1) (Bd.\,I, \S.\,59) is
\weq{5}{D = \Sigma \pm a_{1,1}a_{2,2}\ldots a_{n,n} = A^2.}

The $\alpha_{i,k}$ are then likewise real numbers, which can be determined in infinitely many ways. We assume an arbitrary, but henceforth unalterable, transformation (2), (3).

Let the system of equations (3), solved for the $x$, yield:
\weq{6}{x_i = \beta_{1,i}\xi_1 + \cdots + \beta_{n,i}\xi_n.}

We shall now assign to the variables $x_i$ only integral rational numerical values. It is then clear that the function $f$ can vanish only when all the $x_i$ vanish, and that its value, if we exclude this case, cannot sink below a certain finite limiting value $M$, which we call the \gesperrt{minimum of the form} $f$. For if the value of $f$ sinks below a value $\alpha^2$, then by (2) all the $\xi_i$ must sink below $\alpha$ in absolute value, and by (6), for a sufficiently small $\alpha$, all the $x_i$ would be pressed down below 1 and, since they are integers, would then all have to be equal to zero.

What matters above all is not so much to determine the minimum exactly, as to find for it as small an upper limit as possible, beyond which it cannot go.

We now avail ourselves, with \gesperrt{Minkowski}, of a mode of expression borrowed from a geometry in space of $n$ dimensions, which leads extraordinarily simply to the desired result. For the cases $n=2$, $n=3$ the expressions there employed are completely intuitive and of the very simplest kind. We recommend the reader to make the steps now to be carried out intuitively clear to himself in these two special cases. For greater values of $n$, to be sure, actual intuition is lost. These expressions are then to be regarded as a short notation for certain analytic relations, the analogy with ordinary space still easing their comprehension. There is not the slightest difficulty in replacing all these expressions by purely analytic formulae; but the brevity of expression and the ease of comprehension would suffer greatly by it.

We now designate the variables $\xi_1, \ldots, \xi_n$ as coor-
%CONT>>
